\chapter{\objectivesname}
\label{chap:objetivos}

\section{Objetivo general}

El objetivo general de este Trabajo de Fin de Grado es \textbf{mapear sistem\'aticamente el flujo de trabajo de \textsc{SpecKit} contra un flujo profesional de ingenier\'ia de requisitos basado en normas}, tomando como referencia el IREB CPRE Foundation Level~\cite{cpre}, la norma ISO/IEC/IEEE 29148~\cite{iso29148} y la gu\'ia INCOSE para la redacci\'on de requisitos~\cite{incose}.

\section{Objetivos espec\'ificos}

Para alcanzar el objetivo general, se definen los siguientes objetivos espec\'ificos:

\begin{enumerate}
    \item \textbf{Documentar el flujo SDD de \textsc{SpecKit}}: caracterizar sus comandos, artefactos generados y mecanismos de calidad a partir de sus fuentes oficiales.

    \item \textbf{Definir el flujo normativo de referencia}: establecer un marco de comparaci\'on basado en las actividades IREB, los criterios de calidad de ISO 29148 y las buenas pr\'acticas de INCOSE.

    \item \textbf{Construir un corpus de requisitos reales}: extraer y formalizar requisitos a partir de \emph{changelogs} de seis proyectos \emph{Open Source} representativos de distintos tipos de sistema.

    \item \textbf{Ejecutar el pipeline de \textsc{SpecKit}}: introducir los requisitos formalizados y observar el comportamiento del pipeline SDD, registrando los artefactos generados.

    \item \textbf{Analizar el mapeo}: contrastar los resultados de la ejecuci\'on con el flujo normativo mediante la r\'ubrica SRCI.

    \item \textbf{Identificar patrones y derivar recomendaciones}: extraer conclusiones sobre fortalezas y limitaciones del enfoque SDD desde la perspectiva RE basada en normas.
\end{enumerate}